
\chapter{CONCLUSION AND FUTURE RECOMMENDATIONS}

\ChapterSection{Conclusion}

Krishi Vaidya is a crop health diagnosis project, and it has different component as mobile app, YOLO, VLM, LLM and backend.  Krishi Vaidya runs end to end, and it was verifed thought the field test. The VLM
outputs the structured diagnoses based on the provided image and shows measurable learning, though its
test-set accuracy remains modest after three epochs during training under
computational budget constraints.

The YOLO models had reached F1 0.73 and mAP@50 0.76 across eleven plant-part
classes. Its INT8 quantised variant retains 98 percent of baseline mAP at
1.87$\times$ compression, so based on this we could conclude that on-device deployment is practical. Per-class
recall is lower for small or rare classes like rice panicle and brassica
leaf, as expected since their low sample counts and small sizes in the
training images.

The VLM fine-tuning was trained for three epochs on Qwen2.5-VL-3B with LoRA
adapters and 4-bit quantisation. On the independent test set, 67.7 percent
of predictions failed to parse into a valid label and overall classification
accuracy reached 13.1 percent. These statistics reflect a compute-constrained
baseline: the final evaluation loss fell to 0.0146, mean token accuracy
reached 99.67 percent, and macro precision among parseable predictions stood
at 0.5232. All of the aforementioned results reflect that the model learned the diagnostic task but had not
converged on output within the three-epoch budget. Some steps such as
constrained decoding, stricter prompt templates, and additional training
epochs are the direct levers for closing this gap.

The backend handles authentication, validation, inference calls, and response
formatting. The mobile app covers the farmer-end workflow such as scan capture,
diagnosis display, crop records, and offline operation. 

\ChapterSubsection{Contributions}

\begin{enumerate}[label=\arabic*.]
\item \textbf{Plant-part detection.} YOLOv8 nano was trained on a dataset with multiple crops
and evaluation results were F1 0.73 and mAP@50 0.76 with per-class analysis,
threshold sweeps, and INT8 quantisation retaining 98 percent of baseline mAP.

\item \textbf{Parseability as a VLM metric.} The output-format failures were
evaluated together with classification accuracy, which shows that unparseable
outputs not wrong labels are the main failure mode in structured VLM
diagnosis.

\item \textbf{End-to-end diagnosis workflow.} Krishi Vaidya was field-tested from mobile
image capture through backend validation, VLM inference, and advisory
generation to farmer-facing result display.

\item \textbf{Offline-first mobile design.} The mobile app was designed around an offline-first approach. To achieve this, we have used SQLite persistence, offline scan
queuing, cached advisory fallback, and sync logic built into the mobile app
from initial design rather than added after the fact.

\item \textbf{Constraint reporting.} Documentation of compute-constrained
VLM performance together with evidence of convergence headroom and absence of expert agronomic validation---defining the
boundary between what the current prototype achieves and what extended
training can address.
\end{enumerate}

\ChapterSubsection{Lessons Learned}

\begin{enumerate}[label=\arabic*.]
\item \textbf{Building a model is not the same as building a system.} Training
a detector or fine-tuning a VLM produces an artefact. Wrapping that artefact
in input validation, error handling, uncertainty display, record persistence,
and structured output formatting is a separate effort at least as large
\parencite{sculleyMLTechDebt2015,paleyes2022MLSystemsChallenges}.

    \item \textbf{VLM Performance.} The VLM was trained on limited compute, and it achieved eval loss of 0.0146 and macro precision of 0.523,2 which shows that the model had learned the core task. But it had a parse error rate of 67.7 \% which reflects that the model failed to format the output in the correct order. 

    \item \textbf{Offline architecture.} The Krishi Vaidya project was often used in remote places like crop fields and others. Therefore, the implementation of the mobile app was around the offline-based architecture. \parencite{kleppmannLocalFirst2019}.

\ChapterSection{Future Recommendations}

\ChapterSubsection{Short-term}

\begin{enumerate}[label=\arabic*.]


\item \textbf{VLM output.} To improve the VLM performance, we need to apply constraints to decoding, supervision of the output format during the training phase, a correct prompt template and increased LoRA rank.

\item \textbf{YOLO dataset.} In our current dataset for training the YOLO model, it contains lots of images related to rice panicle, and it has few images of maize leaf; such an imbalance could affect the performance; therefore, the dataset needs to be balanced. 


\end{enumerate}

\ChapterSubsection{Longer-term}

\begin{enumerate}[label=\arabic*.]

\item \textbf{Expert Reviews. } In our current approach of Krishi Vaidya, we are relying on LLM and VLM to produce the diagnosis results; therefore, agricultural experts need to review the diagnosis output. 

\item  \textbf{Farmer uses feasibility. } Krishi Vaidya project needs to provide for lots of farmers and collect feedback on the interface, language, trust of diagnosis, and crop management. 

\item \textbf{ VLM. } We have not tested how the larger-sized VLM would perform, for example,e like 7B or 14B VLM, along with the cost and compute. 

\item \textbf{Continuous learning. } Build infrastructure of user feed back to flag wrong diagnoses and refine components of model.

\end{enumerate}
